\section{Besvarelse av forskningsspørsmålet}
\label{sec:forskningsspm}

Forskningsspørsmålet var:

\begin{quote}
I hvilken grad kan en plugin-basert mellomvarearkitektur standardisere
tilgang til AntMedia Server og NHN Video gjennom ett felles API, samtidig
som backend-spesifikk logikk isoleres fra kjernelogikken i VideoAPI?
\end{quote}

Resultatet viser at en plugin-basert mellomvarearkitektur i stor grad
egner seg for dette formålet, men med noen viktige begrensninger.
VideoAPI klarer å samle de sentrale romoperasjonene bak ett felles API:
oppretting, henting og sletting av rom kan gjøres gjennom samme
endepunkter uavhengig av videomotor. For klienten gir dette et mer
standardisert grensesnitt enn direkte integrasjon mot AntMedia og NHN hver
for seg.

Samtidig viser prosjektet at standardisering ikke betyr at alle
videomotorer kan behandles helt likt. AntMedia og NHN har ulike
domenemodeller. AntMedia er strømorientert, mens NHN er møte- og
bookingorientert. Derfor finnes det operasjoner, særlig start, stopp og
romtoken, som ikke har identisk betydning i begge systemer. VideoAPI løser
dette ved å tilby samme API-flate, men la provider-laget håndtere
forskjellene og returnere tydelige svar når en operasjon ikke støttes.

Den viktigste styrken i løsningen er isoleringen av backend-spesifikk
logikk. Controller-laget og de offentlige modellene forblir uavhengige av
detaljene i AntMedia og NHN. Valideringen med ekstra providere viste også
at nye videomotorer kunne legges til uten endringer i kjernelogikken.
Dette støtter hovedideen bak arkitekturen: nye integrasjoner krever
arbeid, men arbeidet lokaliseres til provider-laget i stedet for å spre
seg gjennom hele API-et.

De ikke-funksjonelle kravene støtter også forskningsspørsmålet. Løsningen
har autentisering, containerisering, miljøbasert konfigurasjon,
helsesjekker, logging og Prometheus/Grafana-metrikker. Dette gjør
arkitekturen mer realistisk som mellomvare i HDOs driftskontekst. Likevel
må resultatet forstås som et proof-of-concept. Løsningen demonstrerer at
arkitekturen fungerer, men den dokumenterer ikke full produksjonsmodenhet.

Samlet er svaret derfor at en plugin-basert mellomvarearkitektur i stor
grad kan standardisere tilgangen til AntMedia og NHN gjennom ett felles
API, så lenge standardiseringen forstås som en felles API-kontrakt og ikke
som full semantisk likhet mellom videomotorene. Backend-spesifikk logikk
kan isoleres fra kjernelogikken, men de underliggende forskjellene må
fortsatt håndteres eksplisitt i provider-laget og dokumenteres for
klientene.